\documentclass[a4paper,12pt]{article}

\usepackage[T2A]{fontenc}
\usepackage[utf8]{inputenc}
\usepackage[russian]{babel}

\usepackage{amsmath,amssymb,amsthm,mathtools}
\usepackage{helvet}
\renewcommand{\familydefault}{\sfdefault}
\usepackage{icomma}
\usepackage{geometry}
\geometry{margin=2cm}
\usepackage{microtype}
\usepackage{tikz}
\usepackage{tkz-euclide}
\usepackage{pgfplots}
\pgfplotsset{compat=1.18}
\usepackage{tabularx,booktabs,longtable}
\usepackage{enumitem}
\usepackage{hyperref}
\usepackage{parskip}
\usepackage{fancyhdr}
\usepackage{setspace}
\usepackage{xcolor}
\usepackage{array}

\definecolor{accent}{HTML}{008080}
\definecolor{darkaccent}{HTML}{004D4D}
\definecolor{textgray}{HTML}{333333}
\definecolor{lightaccent}{HTML}{D9F0F0}

\hypersetup{
	colorlinks=true,
	linkcolor=darkaccent,
	urlcolor=darkaccent
}

\pagestyle{fancy}
\fancyhf{}
\fancyhead[L]{\textcolor{darkaccent}{Чек-лист по алгебре}}
\fancyhead[R]{\textcolor{darkaccent}{ОГЭ}}
\fancyfoot[C]{\thepage}

\onehalfspacing
\setlist[itemize]{leftmargin=2.2em}
\setlist[enumerate]{leftmargin=2.4em}
\setlength{\headheight}{15pt}

\newcommand{\topic}{Графики функций: модуль, алгебраические дроби, выколотые точки, прямые \(y=m\) и \(y=kx\)}
\newcommand{\studentname}{Екатерины Гнедковой}

\begin{document}
\color{textgray}

% ---------------- TITUL ----------------
\thispagestyle{empty}
\vspace*{2.8cm}

\begin{center}
{\Huge\bfseries \textcolor{darkaccent}{Чек-лист}}\\[0.6cm]
{\LARGE \textcolor{accent}{\topic}}\\[1.6cm]

{\Large Лёвин Артём Александрович эксклюзивно для \studentname}\\[1.2cm]

{\large \today}
\end{center}

\vfill

\begin{center}
\begin{tikzpicture}[x=1cm,y=1cm]
\draw[accent,line width=1.1pt]
(0,0) .. controls (2,0.9) and (4,-0.6) .. (6,0.2)
      .. controls (7.5,0.8) and (9,-0.3) .. (11,0.4)
      .. controls (12,0.7) and (13,0.25) .. (14,0.5);
\draw[darkaccent,line width=0.5pt]
(0,-0.22) .. controls (1.5,0.2) and (3.2,-0.4) .. (5,-0.1)
          .. controls (7,0.4) and (9,-0.15) .. (11,0.15)
          .. controls (12.3,0.32) and (13.2,0.05) .. (14,0.18);
\end{tikzpicture}
\end{center}

\newpage

% ---------------- MAIN ----------------
\section*{\textcolor{darkaccent}{Тема занятия}}
На занятии разбирались задачи на построение и чтение графиков функций. Основные типы:
\begin{itemize}
	\item графики с модулем;
	\item графики рациональных выражений после сокращения;
	\item графики с \emph{выколотыми точками};
	\item задачи на количество общих точек с прямыми \(y=m\) и \(y=kx\).
\end{itemize}

\section*{\textcolor{darkaccent}{1. Базовые идеи}}
\begin{itemize}
	\item \textbf{Модуль} раскрывается по определению:
	\[
	|u|=\begin{cases}
	u, & u\ge 0,\\
	-u, & u<0.
	\end{cases}
	\]
	\item \textbf{Алгебраическую дробь} сначала анализируют по ОДЗ, потом раскладывают на множители и только после этого сокращают.
	\item \textbf{После сокращения запрещённые значения не исчезают:} им соответствуют выколотые точки.
	\item \textbf{Прямая \(y=m\)} --- горизонтальная.
	\item \textbf{Прямая \(y=kx\)} проходит через начало координат.
\end{itemize}

\section*{\textcolor{darkaccent}{2. Парабола: что нужно уметь находить быстро}}
Для функции
\[
y=ax^2+bx+c
\]
важно уметь сразу находить:

\begin{itemize}
	\item знак коэффициента \(a\) и направление ветвей;
	\item координату вершины:
	\[
	x_{\text{в}}=-\frac{b}{2a};
	\]
	\item ординату вершины:
	\[
	y_{\text{в}}=f\!\left(x_{\text{в}}\right);
	\]
	\item точки пересечения с осью \(Ox\): решить уравнение \(f(x)=0\);
	\item точку пересечения с осью \(Oy\): подставить \(x=0\).
\end{itemize}

\section*{\textcolor{darkaccent}{3. Алгоритм для графика с модулем}}
Если в формуле есть модуль, работаем так:

\begin{enumerate}
	\item Найдите, при каких \(x\) подмодульное выражение неотрицательно, а при каких отрицательно.
	\item Раскройте модуль отдельно в каждом случае.
	\item Получите график по частям.
	\item Стройте каждую ветвь \textbf{только на своём промежутке}.
	\item На границе случаев проверьте, входит ли точка в график.
\end{enumerate}

\subsection*{Мини-пример}
\[
y=|x|-2=
\begin{cases}
x-2, & x\ge 0,\\
-x-2, & x<0.
\end{cases}
\]
Это две полупрямые с общей вершиной \((0,-2)\).

\section*{\textcolor{darkaccent}{4. Алгоритм для рационального выражения}}
Пусть
\[
y=\frac{P(x)}{Q(x)}.
\]
Тогда:

\begin{enumerate}
	\item Найдите ОДЗ:
	\[
	Q(x)\ne 0.
	\]
	\item Разложите числитель и знаменатель на множители.
	\item Сократите общие множители.
	\item Постройте график упрощённой функции.
	\item Вернитесь к ОДЗ и отметьте выколотые точки.
\end{enumerate}

\subsection*{Мини-пример}
\[
y=\frac{x^2-4}{x-2}=\frac{(x-2)(x+2)}{x-2}=x+2,\qquad x\ne 2.
\]
Значит, график --- прямая \(y=x+2\), но точка \((2,4)\) выколота.

\section*{\textcolor{darkaccent}{5. Выколотая точка}}
Если исходная функция не определена при \(x=a\), а после сокращения формула упростилась, то точку
\[
\bigl(a,f_{\text{упр}}(a)\bigr)
\]
нужно \textbf{выколоть}.

\subsection*{Запомнить}
Если вы забыли выколоть хотя бы одну такую точку, график считается построенным неверно.

\section*{\textcolor{darkaccent}{6. Как читать задачи с параметром}}
\subsection*{Прямая \(y=m\)}
Это горизонтальный уровень. Нужно смотреть, сколько раз этот уровень пересекает график функции.

\begin{itemize}
	\item Для параболы один общий пункт часто получается в вершине.
	\item Если график состоит из нескольких ветвей, надо считать пересечения со \textbf{всем} графиком, а не с отдельной частью.
	\item Выколотая точка не считается общей точкой.
\end{itemize}

\subsection*{Прямая \(y=kx\)}
Эта прямая проходит через начало координат. Если известно, что она проходит через точку \((x_0,y_0)\), где \(x_0\ne 0\), то
\[
k=\frac{y_0}{x_0}.
\]

\section*{\textcolor{darkaccent}{7. Типовые ошибки}}
\begin{itemize}
	\item Неправильно раскрыли модуль.
	\item Построили всю параболу, хотя по условию нужна только часть.
	\item Сократили дробь и забыли ОДЗ.
	\item Не выкололи запрещённую точку.
	\item Ошиблись в вершине из-за знака в формуле
	\[
	x_{\text{в}}=-\frac{b}{2a}.
	\]
	\item Посчитали выколотую точку как реальное пересечение.
	\item Перепутали \(y=m\) и \(y=kx\).
\end{itemize}

\section*{\textcolor{darkaccent}{8. Визуальная памятка: ``сократили, но точку не забыли''}}
Для функции
\[
y=\frac{x^2-4}{x-2}=x+2,\qquad x\ne 2
\]
график выглядит так:

\begin{center}
\begin{tikzpicture}
\begin{axis}[
	width=13cm,
	height=6.8cm,
	axis lines=middle,
	unit vector ratio=1 1,
	samples=200,
	xmin=-4, xmax=5,
	ymin=-3, ymax=7,
	clip=false,
	xlabel={$x$},
	ylabel={$y$},
	xlabel style={below right},
	ylabel style={above left},
	grid=major,
	grid style={dashed, gray!30},
	ticklabel style={font=\small},
	label style={font=\small}
]
\addplot[domain=-4:1.96,accent,thick]{x+2};
\addplot[domain=2.04:5,accent,thick]{x+2};
\addplot[only marks,mark=o,mark size=2.6pt,thick,accent] coordinates {(2,4)};
\node[anchor=west,text=darkaccent,font=\small] at (axis cs:2.2,4.2) {$(2,4)$ --- выколотая точка};
\end{axis}
\end{tikzpicture}
\end{center}

\section*{\textcolor{darkaccent}{9. Визуальная памятка: количество пересечений}}
Рассмотрим функцию
\[
y=\frac{|x^2-4|}{x-2},\qquad x\ne 2.
\]
После раскрытия модуля:
\[
y=\begin{cases}
x+2, & x<-2,\\[1mm]
-(x+2), & -2\le x<2,\\[1mm]
x+2, & x>2.
\end{cases}
\]
Ключевые уровни здесь --- \(y=0\) и \(y=-4\). По рисунку видно, что ровно две общие точки с горизонталью \(y=m\) получаются только при
\[
-4<m<0.
\]

\begin{center}
\begin{tikzpicture}
\begin{axis}[
	width=14cm,
	height=7.2cm,
	axis lines=middle,
	unit vector ratio=1 1,
	samples=200,
	xmin=-6, xmax=6,
	ymin=-6, ymax=7,
	clip=false,
	xlabel={$x$},
	ylabel={$y$},
	xlabel style={below right},
	ylabel style={above left},
	grid=major,
	grid style={dashed, gray!30},
	ticklabel style={font=\small},
	label style={font=\small}
]
\addplot[domain=-6:-2,accent,thick]{x+2};
\addplot[domain=-2:1.98,accent,thick]{-(x+2)};
\addplot[domain=2.02:6,accent,thick]{x+2};

\addplot[darkaccent,densely dashed,domain=-6:6]{0};
\addplot[darkaccent,densely dashed,domain=-6:6]{-4};

\addplot[only marks,mark=*,mark size=1.9pt,accent] coordinates {(-2,0)};
\addplot[only marks,mark=o,mark size=2.6pt,thick,accent] coordinates {(2,4)};

\node[anchor=south west,text=darkaccent,font=\small] at (axis cs:4.3,0.1) {$y=0$};
\node[anchor=south west,text=darkaccent,font=\small] at (axis cs:4.2,-3.9) {$y=-4$};
\node[anchor=west,text=darkaccent,font=\small] at (axis cs:2.25,4.2) {$(2,4)$ выколота};
\end{axis}
\end{tikzpicture}
\end{center}

\section*{\textcolor{darkaccent}{10. Рабочий шаблон оформления}}
\begin{enumerate}
	\item Найти ОДЗ.
	\item Если есть модуль --- разбить на случаи.
	\item Если есть дробь --- разложить на множители и сократить.
	\item Построить упрощённые ветви графика.
	\item Выделить только допустимые части.
	\item Отметить выколотые точки.
	\item Для задач с параметром отдельно объяснить, почему число пересечений именно такое.
\end{enumerate}

\section*{\textcolor{darkaccent}{11. Краткая таблица-ориентир}}
\renewcommand{\arraystretch}{1.35}
\begin{table}[h!]
\centering
\caption{Что означает объект на координатной плоскости}
\begin{tabularx}{\linewidth}{@{}>{\raggedright\arraybackslash}p{0.28\linewidth}X@{}}
\toprule
\textbf{Выражение} & \textbf{Смысл} \\
\midrule
\(y=m\) & Горизонтальная прямая на высоте \(m\) \\
\(y=kx\) & Прямая через начало координат \\
\(x=a\) & Вертикальная прямая \\
\(\dfrac{P(x)}{Q(x)}\), где \(Q(a)=0\) & В точке \(x=a\) функция не определена \\
Сокращение дроби & Возможна выколотая точка \\
\(a>0\) в \(ax^2+bx+c\) & Ветви параболы направлены вверх \\
\(a<0\) в \(ax^2+bx+c\) & Ветви параболы направлены вниз \\
\bottomrule
\end{tabularx}
\end{table}

\newpage

% ---------------- EXERCISES ----------------
\section*{\textcolor{darkaccent}{Тренировочный блок}}

\textbf{Расположение по сложности:} 1--3 --- простые, 4--6 --- средние, 7--9 --- выше средних, 10 --- сложная.

\begin{enumerate}
	\item Постройте график функции
	\[
	y=|x|-2.
	\]

	\item Постройте график функции
	\[
	y=\frac{x^2-4}{x-2}.
	\]

	\item Для функции
	\[
	y=x^2-4x+3
	\]
	найдите вершину и точки пересечения с осями координат.

	\item Постройте график функции
	\[
	y=|x^2-1|.
	\]

	\item Постройте график функции
	\[
	y=\frac{x^2-5x+6}{x-2}.
	\]

	\item Найдите все значения \(m\), при которых прямая
	\[
	y=m
	\]
	имеет с графиком функции
	\[
	y=x^2-2x
	\]
	ровно одну общую точку.

	\item Постройте график функции
	\[
	y=\frac{|x|}{1-|x|}.
	\]

	\item Найдите все значения \(m\), при которых прямая
	\[
	y=m
	\]
	имеет с графиком функции
	\[
	y=\frac{x^2-1}{x-1}
	\]
	ровно одну общую точку.

	\item Найдите все значения \(k\), при которых прямая
	\[
	y=kx
	\]
	проходит через выколотую точку графика функции
	\[
	y=\frac{x^2-5x+6}{x-2}.
	\]

	\item Постройте график функции
	\[
	y=\frac{|x^2-4|}{x-2}
	\]
	и найдите все значения \(m\), при которых прямая \(y=m\) имеет с этим графиком ровно две общие точки.
\end{enumerate}

\end{document}